% ============================================================================
% capa.tex
% ----------------------------------------------------------------------------
% Gera a CAPA do documento (primeira página, sem numeração e sem cabeçalho/
% rodapé padrão).
%
% Usa os dados definidos em:
%   - config.tex  -> caminhos e tamanhos dos logos, cidade
%   - edicao.tex  -> título, subtítulo, ano
%
% Você normalmente NÃO precisa editar este arquivo, apenas os dados que ele
% utiliza (em config.tex e edicao.tex).
% ============================================================================

\begin{titlepage}
    \thispagestyle{empty}   % capa não tem cabeçalho/rodapé nem numeração
    \centering              % centraliza todo o conteúdo da capa

    % -- Logo principal da empresa --------------------------------------------
    % Se o arquivo definido em \logoCapaArquivo (config.tex) não existir,
    % mostra uma caixa com o texto "(logo da empresa aqui)" no lugar —
    % assim o modelo continua compilando mesmo sem as imagens finais.
    \IfFileExists{\logoCapaArquivo}{
        \includegraphics[width=\larguraLogoCapa]{\logoCapaArquivo}
    }{
        \fbox{\rule{0pt}{2cm}\hspace{4cm}\small (logo da empresa aqui)\hspace{4cm}}
    }

    \vspace{0.15\textheight}

    % -- Título e subtítulo do documento ---------------------------------------
    % Vêm de \tituloDocumento e \subtituloDocumento, definidos em edicao.tex.
    {\fontsize{14}{17}\selectfont\bfseries \MakeUppercase{\tituloDocumento}\par}

    \vspace{0.5cm}

    {\fontsize{12}{15}\selectfont \subtituloDocumento\par}

    \vfill

    % -- Logos de parceiros/apoiadores -----------------------------------------
    % Exibidos lado a lado, com opacidade reduzida (\opacidadeLogosParceiras).
    % Cada um também tem um "substituto" em caixa de texto, caso a imagem
    % correspondente não exista no projeto.
    \noindent
    \makebox[\textwidth]{
        \hfill
        \IfFileExists{\logoParceiraUmArquivo}{
            \transparent{\opacidadeLogosParceiras}\includegraphics[height=\alturaLogosParceiras]{\logoParceiraUmArquivo}\transparent{1}
        }{
            \fbox{\rule{0pt}{\alturaLogosParceiras}\hspace{2cm}\small(Logo UnB)\hspace{2cm}}
        }
        \hfill
        \IfFileExists{\logoParceiraDoisArquivo}{
            \transparent{\opacidadeLogosParceiras}\includegraphics[height=\alturaLogosParceiras]{\logoParceiraDoisArquivo}\transparent{1}
        }{
            \fbox{\rule{0pt}{\alturaLogosParceiras}\hspace{2cm}\small(Logo Concentro)\hspace{2cm}}
        }
        \hfill
        \IfFileExists{\logoParceiraTresArquivo}{
            \transparent{\opacidadeLogosParceiras}\includegraphics[height=\alturaLogosParceiras]{\logoParceiraTresArquivo}\transparent{1}
        }{
            \fbox{\rule{0pt}{\alturaLogosParceiras}\hspace{2cm}\small(Logo Brasil Júnior)\hspace{2cm}}
        }
        \hfill
    }

    \vspace{3cm}

    % -- Cidade e ano -----------------------------------------------------------
    {\fontsize{12}{15}\selectfont \cidadeDocumento{}, \anoDocumento\par}

\end{titlepage}

% Reinicia a contagem de páginas: a capa não é contada, então a página
% seguinte (folha de rosto) volta a ser considerada "página 1".
\setcounter{page}{1}